\documentclass[a4paper,landscape,8pt]{extarticle}
\usepackage[landscape,margin=0.7cm,top=0.5cm,bottom=0.7cm]{geometry}
\usepackage[utf8]{inputenc}
\usepackage[T1]{fontenc}
\usepackage[ngerman]{babel}
\usepackage{helvet}
\renewcommand{\familydefault}{\sfdefault}
\usepackage{xcolor}
\usepackage{tikz}
\usetikzlibrary{positioning, shapes.geometric}
\usepackage{enumitem}
\usepackage{multicol}
\usepackage{array}
\usepackage{fancyhdr}
\usepackage{amssymb}

% Farben
\definecolor{macgray}{HTML}{1D1D1F}
\definecolor{finderblue}{HTML}{1A73E8}
\definecolor{dockgray}{HTML}{8E8E93}
\definecolor{menubar}{HTML}{F5F5F7}
\definecolor{trackpadpurple}{HTML}{5856D6}
\definecolor{lightgray}{HTML}{F5F5F5}
\definecolor{bordergray}{HTML}{BBBBBB}
\definecolor{warnred}{HTML}{C62828}
\definecolor{appleblue}{HTML}{007AFF}

\pagestyle{fancy}
\fancyhf{}
\renewcommand{\headrulewidth}{0pt}
\fancyfoot[C]{\footnotesize Seite \thepage{} von 3}

\setlength{\parindent}{0pt}
\setlength{\parskip}{2pt}

% Kompakte Listen
\setlist{nosep, leftmargin=1em, itemsep=0pt, topsep=1pt}

% Tastensymbol
\newcommand{\taste}[1]{\fbox{\small\texttt{#1}}}

\begin{document}

% ============================================================================
% SEITE 1: INFORMATIONSBLATT
% ============================================================================

\noindent{\Large\bfseries Den Mac kennenlernen: Schreibtisch, Menüleiste \& Dock} \hfill
{\footnotesize\textbf{Lernziele:} Desktop verstehen | Menüleiste nutzen | Dock anpassen}

\vspace{1mm}
\hrule
\vspace{1mm}

% Drei-Spalten-Layout
\begin{minipage}[t]{0.32\linewidth}
\begin{center}
\colorbox{macgray!10}{%
\begin{minipage}{0.95\linewidth}
\vspace{2mm}
\begin{center}
{\Large\bfseries\textcolor{macgray}{Der Schreibtisch}}\\[1pt]
{\footnotesize\itshape Dein digitaler Arbeitsplatz}
\end{center}
\vspace{2mm}
\end{minipage}%
}
\end{center}

\vspace{1mm}

\textbf{Was ist der Schreibtisch?} Der Bildschirm, den du siehst, wenn keine Fenster offen sind.

\textbf{Analogie:} Wie dein \textbf{echter Schreibtisch} -- hier legst du Dinge ab und arbeitest.

\textbf{Die Elemente:}
\begin{itemize}
\item \textbf{Menüleiste} = ganz oben (immer da!)
\item \textbf{Schreibtisch} = die große Fläche
\item \textbf{Dock} = die Leiste unten
\item \textbf{Fenster} = offene Programme
\end{itemize}

\textbf{Dateien auf dem Schreibtisch:}
\begin{itemize}
\item Du kannst Dateien hier \textbf{ablegen}
\item Wie Papiere auf dem echten Schreibtisch
\item \textbf{Tipp:} Nicht zu voll machen!
\end{itemize}

\textbf{Schreibtisch zeigen:}\\
Alle Fenster beiseite: \taste{F11} drücken\\
oder mit Trackpad: Daumen + 3 Finger auseinander

\end{minipage}%
\hfill%
\begin{minipage}[t]{0.32\linewidth}
\begin{center}
\colorbox{menubar}{%
\begin{minipage}{0.95\linewidth}
\vspace{2mm}
\begin{center}
{\Large\bfseries\textcolor{macgray}{Die Menüleiste}}\\[1pt]
{\footnotesize\itshape Immer oben, immer wichtig}
\end{center}
\vspace{2mm}
\end{minipage}%
}
\end{center}

\vspace{1mm}

\textbf{Was ist die Menüleiste?} Der \textbf{Streifen ganz oben} am Bildschirm -- ändert sich je nach Programm.

\textbf{Analogie:} Wie das \textbf{Armaturenbrett} im Auto -- zeigt wichtige Infos und Bedienelemente.

\textbf{Links in der Menüleiste:}
\begin{itemize}
\item \textbf{Apfel-Symbol} = Hauptmenü (Neustart, Ausschalten...)
\item \textbf{App-Name} = Menü der aktiven App
\item \textbf{Ablage, Bearbeiten...} = Befehle
\end{itemize}

\textbf{Rechts in der Menüleiste:}
\begin{itemize}
\item \textbf{WLAN-Symbol} = Netzwerk
\item \textbf{Lautsprecher} = Lautstärke
\item \textbf{Batterie} = Akkustand (MacBook)
\item \textbf{Uhrzeit} = Datum und Uhr
\item \textbf{Kontrollzentrum} = Schnelleinstellungen
\end{itemize}

\textbf{Apfel-Menü -- das Wichtigste:}
\begin{itemize}
\item \textbf{Über diesen Mac} = Infos zum Gerät
\item \textbf{Systemeinstellungen} = Alles einstellen
\item \textbf{Neustart / Ausschalten}
\end{itemize}

\end{minipage}%
\hfill%
\begin{minipage}[t]{0.32\linewidth}
\begin{center}
\colorbox{dockgray!15}{%
\begin{minipage}{0.95\linewidth}
\vspace{2mm}
\begin{center}
{\Large\bfseries\textcolor{dockgray}{Das Dock}}\\[1pt]
{\footnotesize\itshape Deine Programmleiste}
\end{center}
\vspace{2mm}
\end{minipage}%
}
\end{center}

\vspace{1mm}

\textbf{Was ist das Dock?} Die \textbf{Leiste mit Symbolen} am unteren Bildschirmrand.

\textbf{Analogie:} Wie eine \textbf{Werkzeugleiste} -- die wichtigsten Programme griffbereit.

\textbf{Dock-Bereiche:}
\begin{itemize}
\item \textbf{Links:} Deine Programme
\item \textbf{Mitte} (nach Strich): Zuletzt benutzt
\item \textbf{Rechts:} Papierkorb + Downloads
\end{itemize}

\textbf{Punkt unter dem Symbol:}\\
$\bullet$ = Programm ist \textbf{geöffnet}

\textbf{Programm zum Dock hinzufügen:}
\begin{enumerate}
\item Programm öffnen (es erscheint im Dock)
\item \textbf{Rechtsklick} auf das Symbol
\item \glqq Im Dock behalten\grqq{} wählen
\end{enumerate}

\textbf{Aus dem Dock entfernen:}\\
Symbol aus dem Dock \textbf{herausziehen} -- es verschwindet mit \glqq Puff\grqq{}.

\end{minipage}

\vspace{1mm}

% Zusammenfassung
\begin{center}
\fcolorbox{bordergray}{lightgray}{%
\begin{minipage}{0.98\linewidth}
\centering
\vspace{1.5mm}
\textbf{Merke:}
\textbf{Menüleiste} = immer oben (ändert sich je nach App) ~$|$~
\textbf{Dock} = immer unten (deine wichtigsten Programme) ~$|$~
\textbf{Schreibtisch} = Arbeitsfläche für Dateien
\vspace{1.5mm}
\end{minipage}%
}
\end{center}

\newpage

% ============================================================================
% SEITE 2: FINDER & TRACKPAD
% ============================================================================

\begin{center}
{\LARGE\bfseries Finder, Trackpad \& wichtige Tastenkombinationen}\\[3pt]
{\normalsize Dateien finden und den Mac effizient bedienen}
\end{center}

\vspace{1mm}
\hrule
\vspace{2mm}

\begin{multicols}{2}

\section*{\textcolor{finderblue}{Der Finder -- Dein Aktenschrank}}

\textbf{Was ist der Finder?} Das Programm zum \textbf{Verwalten deiner Dateien} -- immer aktiv.

\textbf{Symbol:} Das lächelnde Gesicht (blau-weiß) im Dock.

\textbf{Analogie:} Wie ein \textbf{Aktenschrank} mit Ordnern und Dokumenten.

\vspace{2mm}

\textbf{Finder-Fenster öffnen:}
\begin{itemize}
\item Klick auf \textbf{Finder-Symbol} im Dock
\item Oder: \taste{cmd} + \taste{N} (neues Fenster)
\end{itemize}

\textbf{Die Seitenleiste (links im Fenster):}
\begin{itemize}
\item \textbf{Schreibtisch} = Dateien auf dem Desktop
\item \textbf{Dokumente} = Deine Dokumente
\item \textbf{Downloads} = Heruntergeladene Dateien
\item \textbf{Programme} = Alle installierten Apps
\item \textbf{iCloud Drive} = Dateien in der Cloud
\end{itemize}

\textbf{Datei suchen:}
\begin{itemize}
\item Im Finder: \taste{cmd} + \taste{F}
\item Oder: Spotlight öffnen mit \taste{cmd} + \taste{Leertaste}
\end{itemize}

\vspace{2mm}

\textbf{Wichtige Ordner:}

\renewcommand{\arraystretch}{1.3}
\begin{tabular}{@{}p{2.5cm}p{5cm}@{}}
\textbf{Dokumente} & Briefe, Tabellen, Notizen \\
\textbf{Downloads} & Aus dem Internet geladen \\
\textbf{Bilder} & Fotos (nicht Fotos-App!) \\
\textbf{Programme} & Alle installierten Apps \\
\end{tabular}
\renewcommand{\arraystretch}{1}

\vspace{3mm}

\section*{\textcolor{trackpadpurple}{Das Trackpad -- Gesten am Mac}}

\textbf{Was ist das Trackpad?} Die \textbf{berührungsempfindliche Fläche} unter der Tastatur (MacBook) oder das Magic Trackpad.

\textbf{Analogie:} Wie eine \textbf{Maus}, aber mit Fingergesten.

\columnbreak

\textbf{Die wichtigsten Trackpad-Gesten:}

\renewcommand{\arraystretch}{1.4}
\begin{tabular}{@{}p{3.5cm}p{5cm}@{}}
\textbf{Tippen} & Klicken (wie linke Maustaste) \\
\textbf{Mit 2 Fingern tippen} & Rechtsklick (Kontextmenü) \\
\textbf{Mit 2 Fingern scrollen} & Hoch/runter in Dokumenten \\
\textbf{2 Finger auseinander} & Zoomen (vergrößern) \\
\textbf{2 Finger zusammen} & Zoomen (verkleinern) \\
\textbf{3 Finger nach oben} & Alle Fenster zeigen \\
\textbf{3 Finger seitlich} & Zwischen Schreibtischen wechseln \\
\end{tabular}
\renewcommand{\arraystretch}{1}

\vspace{3mm}

\section*{Wichtige Tastenkombinationen}

\taste{cmd} = die Taste mit dem $\clmark$ (Kleeblatt/Apfel)

\vspace{2mm}

\renewcommand{\arraystretch}{1.4}
\begin{tabular}{@{}p{3cm}p{5.5cm}@{}}
\taste{cmd} + \taste{C} & Kopieren \\
\taste{cmd} + \taste{V} & Einsetzen \\
\taste{cmd} + \taste{X} & Ausschneiden \\
\taste{cmd} + \taste{Z} & Rückgängig \\
\taste{cmd} + \taste{A} & Alles auswählen \\
\taste{cmd} + \taste{S} & Speichern \\
\taste{cmd} + \taste{Q} & Programm beenden \\
\taste{cmd} + \taste{W} & Fenster schließen \\
\taste{cmd} + \taste{Leertaste} & Spotlight-Suche \\
\end{tabular}
\renewcommand{\arraystretch}{1}

\vspace{3mm}

\begin{center}
\fcolorbox{appleblue}{appleblue!8}{%
\begin{minipage}{0.9\linewidth}
\vspace{2mm}
\textbf{Tipp: Spotlight-Suche}\\[2pt]
\taste{cmd} + \taste{Leertaste} öffnet die Suche. Hier kannst du \textbf{alles} finden: Programme, Dateien, Kontakte, sogar rechnen!
\vspace{2mm}
\end{minipage}}
\end{center}

\end{multicols}

\newpage

% ============================================================================
% SEITE 3: ÜBUNGEN & CHECKLISTE
% ============================================================================

\begin{center}
{\LARGE\bfseries Übungen \& Checkliste}\\[3pt]
{\normalsize Teste dein Wissen und übe die Mac-Bedienung}
\end{center}

\vspace{1mm}
\hrule
\vspace{2mm}

\begin{multicols}{2}

\textbf{\large Teil A: Begriffe zuordnen}

Was ist was? Schreibe die Nummer.

\vspace{2mm}

\renewcommand{\arraystretch}{1.5}
\begin{tabular}{@{}p{4.5cm}p{0.8cm}p{5cm}@{}}
1. Menüleiste & & A. Die Leiste unten mit Programmen \\
2. Dock & & B. Das Programm für Dateien \\
3. Finder & & C. Die Leiste ganz oben \\
4. Schreibtisch & & D. Schnellsuche für alles \\
5. Spotlight & & E. Die Arbeitsfläche \\
\end{tabular}
\renewcommand{\arraystretch}{1}

\textit{\footnotesize Lösungen: 1-C, 2-A, 3-B, 4-E, 5-D}

\vspace{4mm}

\textbf{\large Teil B: Praktische Übungen}

Führe diese Aktionen auf deinem Mac aus und hake ab:

\vspace{2mm}

\begin{itemize}[label=$\square$, itemsep=3pt]
\item Öffne das \textbf{Apfel-Menü} (oben links)
\item Klicke auf \glqq Über diesen Mac\grqq{}
\item Öffne den \textbf{Finder} (Klick im Dock)
\item Navigiere zu \textbf{Dokumente} in der Seitenleiste
\item Öffne \textbf{Spotlight} mit \taste{cmd} + \taste{Leertaste}
\item Suche nach \glqq Safari\grqq{} und öffne es
\item Füge ein Programm zum \textbf{Dock} hinzu
\item Probiere \textbf{Rechtsklick} mit 2 Fingern auf dem Trackpad
\end{itemize}

\vspace{4mm}

\textbf{\large Teil C: Tastenkombinationen}

Welche Taste(n) drückst du? Verbinde.

\vspace{2mm}

\renewcommand{\arraystretch}{1.5}
\begin{tabular}{@{}p{4cm}p{4cm}@{}}
Kopieren & \taste{cmd} + \taste{Q} \\
Einsetzen & \taste{cmd} + \taste{C} \\
Programm beenden & \taste{cmd} + \taste{V} \\
Suchen & \taste{cmd} + \taste{Leertaste} \\
\end{tabular}
\renewcommand{\arraystretch}{1}

\columnbreak

\textbf{\large Checkliste: Das kann ich jetzt}

\vspace{2mm}

\begin{center}
\fcolorbox{macgray}{macgray!10}{%
\begin{minipage}{0.95\linewidth}
\vspace{2mm}
\textbf{\textcolor{macgray}{Grundlagen}}
\begin{itemize}[label=$\square$, itemsep=2pt]
\item Ich weiß, wo die \textbf{Menüleiste} ist und was das Apfel-Menü enthält
\item Ich weiß, was das \textbf{Dock} ist und kann Programme hinzufügen/entfernen
\item Ich kann den \textbf{Schreibtisch} zeigen
\end{itemize}
\vspace{2mm}
\end{minipage}}
\end{center}

\vspace{2mm}

\begin{center}
\fcolorbox{finderblue}{finderblue!10}{%
\begin{minipage}{0.95\linewidth}
\vspace{2mm}
\textbf{\textcolor{finderblue}{Finder}}
\begin{itemize}[label=$\square$, itemsep=2pt]
\item Ich kann den \textbf{Finder} öffnen
\item Ich kenne die wichtigen Ordner (Dokumente, Downloads...)
\item Ich kann mit \textbf{Spotlight} suchen
\end{itemize}
\vspace{2mm}
\end{minipage}}
\end{center}

\vspace{2mm}

\begin{center}
\fcolorbox{trackpadpurple}{trackpadpurple!10}{%
\begin{minipage}{0.95\linewidth}
\vspace{2mm}
\textbf{\textcolor{trackpadpurple}{Trackpad \& Tasten}}
\begin{itemize}[label=$\square$, itemsep=2pt]
\item Ich kann mit \textbf{2 Fingern} rechtsklicken
\item Ich kann mit \textbf{2 Fingern} scrollen
\item Ich kenne \taste{cmd}+\taste{C} (Kopieren) und \taste{cmd}+\taste{V} (Einsetzen)
\item Ich kann mit \taste{cmd}+\taste{Leertaste} suchen
\end{itemize}
\vspace{2mm}
\end{minipage}}
\end{center}

\vspace{3mm}

\begin{center}
\fcolorbox{bordergray}{white}{%
\begin{minipage}{0.95\linewidth}
\vspace{2mm}
\textbf{Meine Notizen}\\[2pt]
Mein Mac-Modell: \dotfill\\[2mm]
macOS-Version: \dotfill\\[2mm]
Was ich noch üben möchte: \dotfill
\vspace{2mm}
\end{minipage}}
\end{center}

\end{multicols}

\vfill

\begin{center}
\footnotesize\textit{Stand: \today{} -- Der Mac ist anders als iPhone/iPad, aber mit etwas Übung genauso einfach!}
\end{center}

\end{document}
